% --- Άσκηση 38856 (38856.tex) ---
\Askhsh{38856}{4}
\begin{ylh}{2}
    \item 4.1. Ανισώσεις 1ου Βαθμού
    \item 6.3. Η Συνάρτηση $f(x)=ax+\beta$
\end{ylh}
Ο Στέφανος μετακομίζει σε μια νέα πόλη για σπουδές και σκέφτεται να χρησιμοποιεί ποδήλατο για τις καθημερινές του μετακινήσεις. Προβληματίζεται όμως αν τον συμφέρει να αγοράσει καινούριο ποδήλατο ή να νοικιάσει ποδήλατο με το μήνα. Το ποδήλατο που του αρέσει κοστίζει καινούριο 400 ευρώ και υπολογίζει επίσης ότι θα έχει μηνιαία έξοδα για τη συντήρησή του κατά μέσο όρο 10 ευρώ. Ταυτόχρονα ρώτησε σε ένα μαγαζί όπου νοικιάζουν ποδήλατα και η ενοικίαση ποδηλάτου κοστίζει 35 ευρώ τον μήνα, αλλά χωρίς να έχει έξτρα έξοδα συντήρησης.

\begin{alist}
\item Ο Στέφανος θα χρησιμοποιήσει το ποδήλατο για 5 μήνες στις μετακινήσεις του. Πόσο θα πληρώσει συνολικά:


\begin{rlist}
\item
  Αν αγοράσει ποδήλατο;
\monades{3}

\item
  Αν νοικιάσει ποδήλατο;
\monades{3}
\end{rlist}
\item Να γράψετε μία σχέση που να εκφράζει τα συνολικά έξοδα του Στέφανου, αν αγοράσει και χρησιμοποιήσει το ποδήλατο για $x$ μήνες, και μία σχέση για τα έξοδά του, αν νοικιάσει και χρησιμοποιήσει ποδήλατο για $x$ μήνες.
\monades{4}
\item Να υπολογίσετε το χρονικό διάστημα σε μήνες, για το οποίο η αγορά καινούριου ποδηλάτου συμφέρει τον Στέφανο οικονομικά περισσότερο από τη μηνιαία ενοικίαση.
\monades{6}
\item Σύμφωνα με έναν οικονομικό προϋπολογισμό που έκανε ο Στέφανος, τα έξοδά του για τη μετακίνησή του με το ποδήλατο δεν θέλει να ξεπεράσουν τα 600 ευρώ συνολικά τον πρώτο χρόνο των σπουδών του.
\begin{rlist}
\item
  Ποια είναι η μεγαλύτερη χρονική διάρκεια, σε μήνες, που μπορεί να χρησιμοποιήσει το ποδήλατο, αν το αγοράσει, και ποια, αν το νοικιάσει;\monades{6}
\item
  Σε ποια ή ποιες από τις δύο επιλογές (αγορά ή ενοικίαση) τα ετήσια έξοδα για τη μετακίνησή του θα είναι εντός του προϋπολογισμού του;\monades{3}

{\footnotesize Το παραπάνω θέμα αναπτύχθηκε στο πλαίσιο του έργου: «Ανάπτυξη Δοκιμασιών Αξιολόγησης Δεξιοτήτων Εγγραμματισμού στα μαθήματα της Νεοελληνικής Γλώσσας και Λογοτεχνίας, της Άλγεβρας, της Φυσικής και της Χημείας Α' Λυκείου Γενικού Λυκείου» Ανάδοχος: «Ειδικός Λογαριασμός Κονδυλίων Έρευνας (Ε.Λ.Κ.Ε) Πανεπιστημίου Ιωαννίνων» ΑΔΑΜ: 25SYMV016348911 2025-02-20.}
\end{rlist}
\end{alist}
